% TOP_PROBLEM: {"problemId": "problem.bartnik-cosmological-splitting-conjecture", "canonicalTitle": "Bartnik's cosmological splitting conjecture", "releaseRank": 398, "primaryDomain": "mathematical_physics", "primaryDomainLabel": "Mathematical physics", "displayStatus": "open_with_solved_subcases", "releaseStatus": "open_with_solved_subcases"}
% !TeX program = lualatex
% Definition and sources only; this file is not an LLM solution attempt.
\documentclass[11pt]{article}
\usepackage[margin=25mm]{geometry}
\usepackage{fontspec}
\setmainfont{Latin Modern Roman}
\usepackage{amsmath,amssymb,mathtools}
\usepackage{unicode-math}
\setmathfont{Latin Modern Math}
\usepackage{xurl,hyperref}
\hypersetup{colorlinks=true,urlcolor=blue}
\setcounter{secnumdepth}{0}
\setlength{\parindent}{0pt}
\setlength{\parskip}{5pt}
\setlength{\emergencystretch}{3em}
\begin{document}
\section*{\texorpdfstring{Bartnik's cosmological splitting conjecture}{Bartnik's cosmological splitting conjecture}}\label{398}
\subsection{Definitions and mathematical statement}
Use Lorentzian signature $(-,+,\ldots,+)$. A Cauchy hypersurface meets every inextendible timelike curve exactly once, and global hyperbolicity supplies a well-behaved causal evolution. Let $(M,g)$ be a connected smooth time-oriented boundaryless spacetime of dimension $n\ge2$, globally hyperbolic with a compact Cauchy hypersurface, and complete along timelike geodesics in both time directions. Assume timelike convergence:
\[
 \operatorname{Ric}_g(v,v)\ge0\qquad\text{for every timelike }v.
\]
Bartnik's conjecture asks for a global isometry $(M,g)\cong({\mathbb{R}}\times S,-{\,\mathrm{d}} t^2+h)$ with $(S,h)$ compact Riemannian. The splitting is metric and global, not just the topological product supplied by global hyperbolicity. No timelike line or stationary symmetry is added as an assumption; obtaining the splitting without those extra hypotheses is the target.
\par\smallskip\textit{Formulation and concept references:} \href{https://arxiv.org/html/1810.03183v2}{[S1]}.

\subsection{Short English statement}
Must every timelike-complete spatially compact spacetime satisfying the timelike Ricci condition be a static metric product of time with a compact space?
\subsection{Sources}
\begin{itemize}
\item[S1]Costa e Silva–Flores–Herrera: Conformal symmetries and Bartnik splitting.
\url{https://arxiv.org/html/1810.03183v2}.
\textit{Formulation source.}
\end{itemize}
\end{document}
